\section{Best Time to Buy and Sell Stock with Cooldown}
\label{sec:dp-stock-5}

\problemstatement{Given $\textit{price}$, find the maximum profit from
unlimited transactions, with one restriction: after selling, you must
wait at least one full day (a \textbf{cooldown}) before buying again.}

\begin{patternbox}
\emph{``Cooldown after selling''} modifies exactly one transition in
Section~\ref{sec:dp-stocks-intro}'s base template: selling no longer
leads to $f(i{+}1, 0)$ -- it leads to $f(i{+}2, 0)$, skipping the
mandatory cooldown day entirely.
\end{patternbox}

\subsection*{Deriving the recurrence}

\[
  f(i, \textit{holding}) =
  \begin{cases}
    \max\big(-\textit{price}[i] + f(i{+}1, 1),\ f(i{+}1, 0)\big) &
    \text{not holding}, \\
    \max\big(\textit{price}[i] + f(i{+}2, 0),\ f(i{+}1, 1)\big) &
    \text{holding.}
  \end{cases}
\]
Only the sell clause changed: $f(i{+}1,0) \to f(i{+}2,0)$.

\subsection*{Approach 1: Recursion}

\begin{lstlisting}
#include <bits/stdc++.h>
using namespace std;

int stock5Recursive(int i, int holding, vector<int>& price) {
    int n = price.size();
    if (i >= n) return 0;

    if (holding) {
        int sell = price[i] + stock5Recursive(i + 2, 0, price);
        int hold = stock5Recursive(i + 1, 1, price);
        return max(sell, hold);
    }
    int buy = -price[i] + stock5Recursive(i + 1, 1, price);
    int skip = stock5Recursive(i + 1, 0, price);
    return max(buy, skip);
}

int main() {
    vector<int> price = {1, 2, 3, 0, 2};
    cout << stock5Recursive(0, 0, price) << endl; // 3
    return 0;
}
\end{lstlisting}

\begin{complexitybox}[Approach 1 Complexity]
\textbf{Time.} $O(2^n)$.
\textbf{Space.} $O(n)$ recursion stack.
\end{complexitybox}

\subsection*{Approach 2: Memoization}

\begin{lstlisting}
#include <bits/stdc++.h>
using namespace std;

int stock5Memo(int i, int holding, vector<int>& price, vector<vector<int>>& memo) {
    int n = price.size();
    if (i >= n) return 0;
    if (memo[i][holding] != -1) return memo[i][holding];

    int result;
    if (holding) {
        result = max(price[i] + stock5Memo(i + 2, 0, price, memo),
                     stock5Memo(i + 1, 1, price, memo));
    } else {
        result = max(-price[i] + stock5Memo(i + 1, 1, price, memo),
                     stock5Memo(i + 1, 0, price, memo));
    }
    return memo[i][holding] = result;
}

int main() {
    vector<int> price = {1, 2, 3, 0, 2};
    int n = price.size();
    vector<vector<int>> memo(n, vector<int>(2, -1));
    cout << stock5Memo(0, 0, price, memo) << endl; // 3
    return 0;
}
\end{lstlisting}

\begin{complexitybox}[Approach 2 Complexity]
\textbf{Time.} $O(n)$.
\textbf{Space.} $O(n)$ memo $+$ $O(n)$ recursion stack.
\end{complexitybox}

\subsection*{Approach 3: Tabulation}

\begin{ideabox}[Watch the Index Arithmetic at the Boundary]
Because the sell clause now reaches two days ahead, the tabulation
table needs \textbf{two} extra rows beyond $n$ as a buffer (indices $n$
and $n+1$), both initialized to $0$, rather than just one extra row as
in every earlier section.
\end{ideabox}

\begin{lstlisting}
#include <bits/stdc++.h>
using namespace std;

int stock5Tabulation(vector<int>& price) {
    int n = price.size();
    vector<vector<int>> dp(n + 2, vector<int>(2, 0));

    for (int i = n - 1; i >= 0; i--) {
        dp[i][1] = max(price[i] + dp[i + 2][0], dp[i + 1][1]);
        dp[i][0] = max(-price[i] + dp[i + 1][1], dp[i + 1][0]);
    }
    return dp[0][0];
}

int main() {
    vector<int> price = {1, 2, 3, 0, 2};
    cout << stock5Tabulation(price) << endl; // 3
    return 0;
}
\end{lstlisting}

\subsection*{Dry run}

\dryrun{Take $\textit{price} = [1,2,3,0,2]$ (days $0$--$4$).}

Filling the table backward gives $\textit{dp}[0][0]=3$, realized by two
separate transactions: buy on day $0$ (price $1$), sell on day $1$
(price $2$, profit $1$); cooldown on day $2$; buy on day $3$ (price
$0$), sell on day $4$ (price $2$, profit $2$). Total profit
$1+2=3$. Note that selling on day $1$ forces the next buy to wait until
day $3$ (day $2$ is the mandatory cooldown), which the recurrence
enforces automatically via its $i{+}2$ jump on every sell.

\begin{complexitybox}[Approach 3 Complexity]
\textbf{Time.} $O(n)$.
\textbf{Space.} $O(n)$ table, no recursion stack.
\end{complexitybox}

\subsection*{Approach 4: Space Optimization}

\begin{ideabox}[Now Three Rolling Rows Are Needed, Not Two]
Because the sell clause reads $\textit{dp}[i+2][0]$, the space-optimized
version must keep track of \emph{two} steps of history for the
not-holding value, not just one -- three rolling variables (or two small
arrays) instead of the usual two.
\end{ideabox}

\begin{lstlisting}
#include <bits/stdc++.h>
using namespace std;

int stock5SpaceOptimized(vector<int>& price) {
    int n = price.size();
    int holdNext1 = 0, notHoldNext1 = 0; // i+1
    int notHoldNext2 = 0;                // i+2

    for (int i = n - 1; i >= 0; i--) {
        int curHold = max(price[i] + notHoldNext2, holdNext1);
        int curNotHold = max(-price[i] + holdNext1, notHoldNext1);

        notHoldNext2 = notHoldNext1;
        notHoldNext1 = curNotHold;
        holdNext1 = curHold;
    }
    return notHoldNext1;
}

int main() {
    vector<int> price = {1, 2, 3, 0, 2};
    cout << stock5SpaceOptimized(price) << endl; // 3
    return 0;
}
\end{lstlisting}

\begin{complexitybox}[Approach 4 Complexity]
\textbf{Time.} $O(n)$.
\textbf{Space.} $O(1)$ -- three scalars instead of two.
\end{complexitybox}

\begin{takeawaybox}
A cooldown (or any fixed-length mandatory delay) changes a recurrence's
\emph{lookback/lookahead distance} for one specific transition, which
ripples into two places: the tabulation table needs extra boundary rows
sized to the longest jump, and the space-optimized version needs to keep
that many extra rows of rolling history, not just the usual one.
\end{takeawaybox}
